%!TEX root = ../main.tex
% =====================================================================
% MATERIALS AND METHODS
% "Know-When-to-Trust Crop-Row Guidance: Conformal Heading Intervals
%  and a Risk-Controlled Steering Policy on a Stock Edge Detector"
%
% Honest positioning: the contribution is the TRUST LAYER (conformal
% heading intervals + risk-controlled steering policy) on a STOCK
% detector. The deployed line fit is robust Huber-IRLS. The DFL-variance
% precision-weighting / calibration thesis (old "CALM-Row") is
% empirically dead (rho = -0.13) and appears ONLY as a negative control
% (E0/E3). No new architecture, no new NPU operator as a headline claim.
%
% Equations / protocol match conformal.py, eval_guidance.py,
% eval_conformal.py. No experimental numbers are stated here.
% =====================================================================

\section{Materials and methods}
\label{sec:method}

\noindent We do \emph{not} propose a new backbone, neck, or detection head, and we do not propose a
new line estimator. The detector is a stock, unmodified, edge-deployable lightweight YOLOv8n~\cite{ct4}
(with YOLOv10n~\cite{yolov10} used as a secondary detector where noted), and the guidance line is fit
with a standard robust Huber iteratively-reweighted-least-squares (IRLS) procedure. The contribution of
this paper is a \emph{trust layer} that sits on top of this stock readout: a label-free nonconformity
score, a split-conformal calibration that turns it into finite-sample-valid heading prediction
intervals, and a risk-controlled abstain/slow/proceed steering policy. We position the conformal layer
as a \emph{domain instantiation} of conformal prediction~\cite{vovk2005,lei2018,angelopoulos2023} on a
detector and explicitly contrast it against prior conformal-on-detector
work~\cite{confdet1,confdet2}; we do not claim primacy.

The trust layer runs entirely on the host CPU over post-NMS detections, so it adds no operators to the
detector graph exported to the NPU/accelerator. We state this once, plainly, as a deployment property
rather than a selling point: because the readout is CPU post-NMS arithmetic over quantities the detector
already produces, the exported INT8/RKNN~\cite{rknntoolkit} or TensorRT~\cite{tensorrt} graph remains a
vanilla YOLO. Figure~\ref{fig:pipeline} summarises the pipeline. Because no experiments have been run
yet, every quantitative cell in this paper is a placeholder (``\texttt{-{}-}'' in tables, ``[TODO]'' in
prose); the table and figure scaffolds in Section~\ref{sec:results} are wrapped so that the document
compiles before the numbers exist.

%!TEX root = ../main.tex
% Know-When-to-Trust pipeline overview (TikZ; no external assets).
% Honest proposal: stock detector + robust Huber-IRLS fit (NPU graph is a
% vanilla YOLO), then a CPU post-NMS TRUST LAYER (LOO score -> split-conformal
% interval -> abstain/slow/proceed policy). The DFL-variance precision-weighting
% thesis is dead and does NOT appear here.
\begin{figure}[!t]
  \centering
  \resizebox{\linewidth}{!}{%
  \begin{tikzpicture}[
    node distance=7mm and 8mm,
    box/.style={draw, rounded corners=2pt, align=center, font=\small,
                inner sep=4pt, minimum height=13mm, text width=24mm},
    npu/.style={box, fill=blue!8, draw=blue!50},
    cpu/.style={box, fill=orange!12, draw=orange!60!black},
    io/.style={box, fill=black!5, text width=18mm},
    flow/.style={-{Latex[length=2.4mm]}, thick},
  ]
    % ---- main forward flow (top row) ----
    \node[io]  (img)                   {Input frame};
    \node[npu, right=of img]  (det)    {Stock YOLO\\backbone\,+\,neck\\+\,head\\\scriptsize\cite{ct4,yolov10}};
    \node[npu, right=of det]  (box)    {Post-NMS\\row boxes};
    \node[cpu, right=of box]  (sel)    {Select\\central row};
    \node[cpu, right=of sel]  (fit)    {Robust Huber\\IRLS fit\\$x{=}ay{+}b$\\$\theta{=}\arctan a$};

    % ---- second row: trust layer ----
    \node[cpu, below=14mm of fit, text width=26mm] (score)
        {LOO heading\\dispersion $s$};
    \node[cpu, left=10mm of score, text width=30mm] (conf)
        {Split-conformal\\interval $[\theta{-}h,\theta{+}h]$\\$h{=}\hat q\,s$};
    \node[io, left=10mm of conf, text width=26mm] (pol)
        {Policy:\\proceed / slow /\\abstain};

    \draw[flow] (img)--(det);
    \draw[flow] (det)--(box);
    \draw[flow] (box)--(sel);
    \draw[flow] (sel)--(fit);
    \draw[flow] (fit)--(score);
    \draw[flow] (score)--(conf);
    \draw[flow] (conf)--(pol);

    % ---- grouping boxes (NPU vs CPU trust layer) ----
    \begin{scope}[on background layer]
      \node[draw=blue!50, dashed, rounded corners, fit=(det)(box),
            inner sep=5mm, label={[blue!50!black,font=\scriptsize]above:On-device NPU graph (INT8 RKNN / TensorRT)\,---\,vanilla YOLO}] {};
      \node[draw=orange!60!black, dashed, rounded corners, fit=(sel)(fit)(score)(conf)(pol),
            inner sep=5mm, label={[orange!60!black,font=\scriptsize]below:CPU, post-NMS\,---\,trust layer (no added NPU operators)}] {};
    \end{scope}
  \end{tikzpicture}}
  \caption{Know-when-to-trust crop-row guidance pipeline. A \emph{stock} YOLO detector (blue) emits
  post-NMS crop-row boxes; on the host CPU (orange) the central navigation row is selected, a robust
  Huber-IRLS line $x=ay+b$ is fit and its heading $\theta=\arctan a$ read out. The contribution is the
  trust layer: a label-free leave-one-out heading-dispersion score $s$ feeds a split-conformal
  calibration that produces a finite-sample-valid heading interval of half-width $h=\hat q\,s$, which in
  turn drives a risk-controlled proceed/slow/abstain steering policy. Everything after NMS is CPU
  arithmetic over quantities the detector already outputs, so the exported edge graph remains a vanilla
  YOLO. The detector's native DFL box-distribution variance is shown only as a non-informative negative
  control and is not used here.}
  \label{fig:pipeline}
\end{figure}


% =====================================================================
\subsection{Notation}
\label{sec:method:notation}
% =====================================================================

\noindent We use the following symbols throughout (single source of truth; no synonyms). The guidance
line is parameterised as
\begin{equation}
\label{eq:line}
x = a\,y + b ,
\end{equation}
i.e.\ image column $x$ as a function of image row $y$. This orientation is deliberate: crop rows viewed
head-on are near-vertical, so a conventional $y=mx+c$ fit becomes ill-conditioned as the slope diverges;
expressing $x$ as a function of $y$ removes this vertical-line degeneracy~\cite{sogaard}. The heading is
\begin{equation}
\label{eq:heading}
\theta = \arctan(a),
\end{equation}
reported in degrees from the vertical as the controller-independent navigation quantity, and the
cross-track offset $e_{ct}$ follows from $(a,b)$ evaluated at a look-ahead row (px; cm only under a
calibrated homography, Section~\ref{sec:setup-homography}). The detected row centroids are
$\{(x_i,y_i)\}_{i=1}^{n}$ (post-NMS box centres on the selected row). The robust line fit uses
Huber-IRLS weights $w_i$ from the Huber influence function. The nonconformity (trust) score is
$s$ = leave-one-out (LOO) heading dispersion in degrees. The held-out calibration scores are
$\{s_1,\dots,s_m\}$, split \emph{by base image}; $\hat{q}$ is the split-conformal quantile, the
$\lceil (m+1)(1-\alpha)\rceil$-th smallest of the calibration nonconformity values; the interval
half-width is $h=\hat{q}\,s$; the heading interval is $[\theta-h,\theta+h]$; $\alpha$ is the nominal
miscoverage (we report $\alpha\in\{0.10,0.05\}$); the target coverage is $1-\alpha$ and the realised
coverage is the empirical fraction of test frames whose true heading lies in the interval. The policy
thresholds are $\tau_p$ (proceed) and $\tau_s$ (slow), with abstain otherwise. The symbol
$\sigma^2_{\mathrm{DFL}}$ denotes the detector's distribution-focal-loss (DFL)~\cite{gfl} box-edge
distribution variance; it appears \emph{only} as the non-informative negative-control score in E0/E3 and
is never part of the deployed method.

% =====================================================================
\subsection{Datasets and ground truth}
\label{sec:setup-data}
% =====================================================================

\noindent\textbf{Detection labels from line ground truth.} Neither public benchmark provides bounding
boxes. CRDLD~\cite{cth1} ground truth is a binary mask of crop-row \emph{centrelines}; CRBD~\cite{crbd}
provides parametric (per-scanline) row annotations. As the detector is a stock DFL box detector, we
synthesise detection labels from the line ground truth by placing boxes on \emph{all} visible crop rows
(a single \texttt{crop\_row} segment class) at successive image rows, perspective-scaled in size.
Labelling all rows rather than only the central one avoids contradictory supervision (the rows are
visually identical), which yields a consistent, learnable target; the central navigation row is then
selected from the detections at \emph{evaluation} time (Section~\ref{sec:method:pipeline}). The box
labels are clean, so the detector still learns image difficulty rather than label noise.

\noindent\textbf{In-domain training and test: CRDLD.} The Crop Row Detection dataset of De~Silva
\textit{et~al.}~\cite{cth1} ($512\times512$; train/validation/test $=1250/250/430$) trains the detector
and serves as the in-domain test set on which the conformal coverage guarantee is evaluated under
exchangeability (E1--E4). It carries no per-field metadata, so a within-CRDLD leave-one-field-out
protocol is unavailable.

\noindent\textbf{Cross-dataset transfer: CRBD.} The Crop Row Benchmark Dataset~\cite{crbd}
($320\times240$; a different camera, crop and resolution; no training split) is a held-out
cross-dataset transfer test: the CRDLD-trained detector is evaluated on CRBD without retraining. Because
the two distributions differ, CRBD is used as a coverage-\emph{degradation} study, \emph{not} as a
setting in which the finite-sample coverage guarantee holds (Section~\ref{sec:method:conformal}).

\noindent\textbf{Field video and labelling protocol (C3).} We additionally collect and release a
diverse public crop-row \emph{field-video} dataset of real maize (forward-facing video; clips spanning
field conditions), with a hand-labelled evaluation subset of at least 100 frames. For each labelled
frame an annotator marks the central navigation row, from which the central-row line is fit in the same
$x=a\,y+b$ parameterisation~\eqref{eq:line}; inter-annotator agreement is reported in degrees
(Section~\ref{sec:setup-gtline}). The field video is uncalibrated, so its errors are reported in pixels
and degrees only, never cm; it is used as a second cross-distribution coverage-degradation probe (E2,
E5) and as the ordered frame sequence for the open-loop policy replay (E6). Frames from the same clip
are kept together when splitting, to respect the exchangeability assumption as far as the data allow.

\noindent The CRDLD train/validation/test splits are disjoint, CRBD is an entirely separate dataset, and
the field video is recorded independently, so memorisation cannot inflate any transfer or coverage
result. Table~\ref{tab:setup-datasets} summarises the role of each dataset. A Data and Code Availability
statement accompanies the manuscript.

\begin{table}[!ht]
\centering
\caption{Datasets and their role. Boxes are synthesised from the line ground truth
(Section~\ref{sec:setup-data}); the evaluation ground-truth line is annotation/mask-derived and
detector-independent. cm units are reported only where the camera is calibrated. \emph{Results to be
completed.}}
\label{tab:setup-datasets}
\resizebox{\columnwidth}{!}{\begin{tabular}{lllc}
\hline
\textbf{Dataset} & \textbf{Role} & \textbf{GT guidance line} & \textbf{cm scale?} \\
\hline
CRDLD~\cite{cth1} ($512^2$; 1250/250/430) & Train detector $+$ in-domain coverage & line-mask $\rightarrow$ central row & calibrated \\
CRBD~\cite{crbd} ($320{\times}240$) & Cross-dataset coverage-degradation & annotation $\rightarrow$ central row & px/deg \\
Field video (maize; $\geq$100 labelled) & Coverage-degradation $+$ E6 replay & hand-labelled central row & px/deg \\
\hline
\end{tabular}}
\end{table}

% =====================================================================
\subsection{Pixel-to-centimetre calibration}
\label{sec:setup-homography}
% =====================================================================

\noindent Because ``cm'' is meaningless without a metric anchor, we report centimetres only where the
camera is calibrated (the robot platform and the CRDLD homography), and pixels/degrees everywhere else.
Where camera intrinsics and extrinsics (or a checkerboard reference) are available, we recover the
image-to-ground homography $H$ under an explicit \emph{flat-ground assumption} and report it; the
projection $H$ is applied to line samples to obtain cross-track in cm. Otherwise --- including the field
video --- we report errors in pixels and in pixels normalised by image width, and do not invent a
centimetre value. Flat ground is stated as an explicit limitation with an error budget: a camera-pitch
perturbation propagates to a look-ahead-distance error that grows with range, so we tabulate the induced
centimetre error at the look-ahead row as a function of pitch, making the conditions under which the
centimetre numbers are trustworthy auditable.

% =====================================================================
\subsection{Frozen ground-truth guidance-line extraction}
\label{sec:setup-gtline}
% =====================================================================

\noindent The ground-truth (GT) guidance line used for \emph{evaluation} is derived from the
segmentation mask / annotation and is therefore detector-independent. The extraction is frozen as
follows (\texttt{gt\_line\_from\_mask}). For each crop-row mask we compute, per image row $y$, the
median column $x$ of foreground pixels, then fit a robust Huber-IRLS regression to those per-row
medians, selecting as the navigation row the one nearest the image-centre-bottom; gap and occlusion
handling are part of the frozen procedure. The resulting line is expressed in the same $x=a\,y+b$
parameterisation~\eqref{eq:line}, with heading $\theta=\arctan(a)$~\eqref{eq:heading}. To validate the
auto-extracted GT we cross-check it against a human-annotated subset of at least 100 frames and report
inter-annotator agreement in degrees (and cm where the homography permits), in
Table~\ref{tab:setup-gtline}. The GT line is mask/annotation-derived, not a least-squares fit through GT
box centroids, so it is independent of the box-synthesis step of Section~\ref{sec:setup-data}.

\begin{table}[!ht]
\centering
\caption{Frozen GT-line extraction: agreement between the automatic mask-derived line and human
annotation on the cross-check subset ($\geq$100 frames). cm reported only on the calibrated platform.
\emph{Results to be completed.}}
\label{tab:setup-gtline}
\begin{tabular}{lcc}
\hline
\textbf{Quantity} & \textbf{Heading ($^\circ$)} & \textbf{Cross-track (cm)} \\
\hline
Auto vs.\ human (median $|\Delta|$) & -- & -- \\
Inter-annotator (median $|\Delta|$) & -- & -- \\
\hline
\end{tabular}
\end{table}

% =====================================================================
\subsection{The guidance pipeline: detect, select, robust fit}
\label{sec:method:pipeline}
% =====================================================================

\noindent\textbf{Detect.} Given a forward-facing field image, the stock detector returns crop-row
segment boxes after non-maximum suppression (NMS). For box $i$ we take its centroid $(x_i,y_i)$ in
pixels (\texttt{detect\_with\_logits} $\rightarrow$ \texttt{guidance\_line} in \texttt{eval\_guidance.py}).
The detector fires on \emph{all} rows by design.

\noindent\textbf{Select the central row.} A single guidance line is fit through the centroids of the
\emph{central} navigation row. We seed at the bottom-centre detection, estimate a local slope from the
bottom-most boxes near the seed, and keep detections within a fixed image-fraction corridor of the
\emph{predicted} central-row line $a_{\mathrm{seed}}\,y + b_{\mathrm{seed}}$ (\texttt{select\_central}).
Selecting against the predicted line rather than walking from box to box prevents the reference from
drifting onto an adjacent row, which would leak wrong-row outliers into the fit --- the dominant failure
mode this paper's trust score is designed to catch.

\noindent\textbf{Robust Huber-IRLS line fit (deployed).} Given the $n$ selected centroids, the deployed
guidance line solves a Huber robust regression of $x$ on $y$,
\begin{equation}
\label{eq:irls_obj}
(\hat a,\hat b) = \arg\min_{a,b}\; \sum_{i=1}^{n} \rho_\delta\!\bigl(x_i - a\,y_i - b\bigr),
\qquad
\rho_\delta(r) =
\begin{cases}
\tfrac{1}{2}r^2, & |r|\le\delta,\\[2pt]
\delta\bigl(|r|-\tfrac{1}{2}\delta\bigr), & |r|>\delta,
\end{cases}
\end{equation}
solved by iteratively reweighted least squares (IRLS): at each iteration the residuals
$r_i = x_i-\hat a\,y_i-\hat b$ are rescaled by a robust scale $\hat\sigma = 1.4826\,\mathrm{median}_i|r_i|$
and the Huber influence weights
\begin{equation}
\label{eq:huber_w}
w_i = \min\!\left(1,\ \frac{\delta\,\hat\sigma}{|r_i|}\right),
\qquad \delta = 1.345,
\end{equation}
are applied in a weighted least-squares solve of $x=a\,y+b$; a small per-diagonal ridge keeps the
normal equations well-conditioned (the intercept and slope entries differ by orders of magnitude). The
weights $w_i$ are the \emph{Huber} weights of Equation~\eqref{eq:huber_w} --- robustness to wrong-row
contamination --- and are \emph{not} DFL precision weights. We verified that the detector's native DFL
box-distribution variance $\sigma^2_{\mathrm{DFL}}$ is non-informative for heading error in this setting
(Spearman $\rho=-0.13$; reported as the E0 negative control, Section~\ref{sec:res-e0-trustinfo}), so
robustness, not variance weighting, is what works. This robust readout matches a RANSAC baseline and far
exceeds equal-weight least squares (C1, E1).

\noindent\textbf{Baseline and ablation readouts.} The harness exposes six readouts sharing one stock
checkpoint, so they are compared paired per frame: \texttt{equalLS} (non-robust equal-weight LS, the
detect-then-fit recipe of~\cite{diao_yolox}); \texttt{ransac} (robust consensus LS, a strong baseline);
\texttt{irls} (the deployed Huber-IRLS, equal data weights); \texttt{calm} (the dead DFL-precision
$\times$ Huber-IRLS variant, retained only as a negative control --- it is \emph{not} the
contribution); \texttt{qual} (detector-confidence-squared weights); and \texttt{oracle} (weights from
the true residual to the GT line, an upper bound). Heading error is
$|\arctan(\hat a)-\arctan(a_{\mathrm{gt}})|$ and line-fit RMSE is sampled over $y$ \emph{including} the
far look-ahead row (\texttt{heading\_err\_deg}, \texttt{line\_rmse\_px}).

% =====================================================================
\subsection{Label-free trust score: leave-one-out heading dispersion}
\label{sec:method:score}
% =====================================================================

\noindent The trust score must be \emph{label-free} (computable at run time without ground truth) and
\emph{matched to the failure mode} (wrong-row contamination, which inflates heading error). We use
leave-one-out (LOO) heading dispersion (\texttt{loo\_heading\_dispersion}). For each of the $n$ central
centroids, we refit the guidance line with that centroid removed and record the resulting heading
$\theta^{(-i)}=\arctan(\hat a^{(-i)})$; the nonconformity score is the dispersion of these LOO headings,
\begin{equation}
\label{eq:loo_score}
s = \operatorname{std}_{i=1,\dots,n}\!\bigl(\theta^{(-i)}\bigr)\ \ (\text{deg}).
\end{equation}
A high $s$ means the heading hinges on a single detection --- a wrong-row outlier, or too few /
ill-conditioned points --- and should not be trusted; a low $s$ means the heading is geometrically
stable. This score directly targets the measured failure mode, unlike $\sigma^2_{\mathrm{DFL}}$. As
candidate alternatives for the sharpness ablation (E3) the harness also logs: the robust residual scale
of the fit, $1-$inlier-fraction, $1-$mean detection confidence, the mean DFL centroid sigma
($\sigma_{\mathrm{DFL}}$, the dead control), and $1/\sqrt{n}$ (\texttt{s\_resid}, \texttt{s\_inv\_inlier},
\texttt{s\_conf}, \texttt{s\_dfl}, \texttt{s\_ninv}). The mean detection confidence is also retained as
the score for the fixed-confidence baseline gate (E4).

% =====================================================================
\subsection{Split-conformal heading intervals}
\label{sec:method:conformal}
% =====================================================================

\noindent We turn the score $s$ into a finite-sample-valid heading prediction interval via
\emph{normalised (score-scaled) split conformal}~\cite{vovk2005,lei2018,romano2019cqr} (\texttt{conformal.py}).
On a held-out calibration set of $m$ frames we form the normalised residuals
\begin{equation}
\label{eq:conf_R}
R_j = \frac{|\theta_j - \theta_j^{\mathrm{gt}}|}{s_j}, \qquad j=1,\dots,m,
\end{equation}
and take the split-conformal quantile
\begin{equation}
\label{eq:conf_qhat}
\hat q = R_{(k)},\qquad k = \bigl\lceil (m+1)(1-\alpha)\bigr\rceil,
\end{equation}
the $k$-th smallest of $\{R_j\}$ (\texttt{calibrate}). For a test frame with heading $\theta$ and score
$s$ the interval is $[\theta-h,\theta+h]$ with half-width
\begin{equation}
\label{eq:conf_hw}
h = \hat q\,s
\end{equation}
(\texttt{half\_width}). The constant-width conformal baseline sets $s\equiv 1$.

\noindent\textbf{Coverage statement and exchangeability (honest).} Under \emph{exchangeability of the
calibration and test scores within one distribution}, split conformal guarantees the finite-sample
bound $\Pr\!\bigl(|\theta-\theta^{\mathrm{gt}}|\le \hat q\,s\bigr)\ge 1-\alpha$, i.e.\ realised coverage
$\ge 1-\alpha$ in expectation over the calibration draw~\cite{vovk2005,lei2018}. This guarantee holds
on the in-domain CRDLD test (E2), where we split \emph{by base image} to keep frames from the same scene
out of both calibration and test. It does \emph{not} transfer to CRBD or the field video, whose
distributions differ; for those we report realised coverage as a \emph{coverage-degradation study},
explicitly not a guarantee, and we state this every time coverage is claimed. Realised coverage is
reported with Clopper--Pearson (and Wilson) confidence intervals (\texttt{clopper\_pearson},
\texttt{wilson\_ci}), marginally and stratified by difficulty.

% =====================================================================
\subsection{Risk-controlled steering policy and open-loop replay}
\label{sec:method:policy}
% =====================================================================

\noindent\textbf{Abstain/slow/proceed.} The certified half-width $h=\hat q\,s$ drives a three-way
steering policy: \emph{proceed} if $h\le\tau_p$, \emph{slow} if $\tau_p<h\le\tau_s$, and \emph{abstain}
(stop / hold the last trusted heading) if $h>\tau_s$. Sweeping the threshold traces a selective-risk
curve --- heading error on \emph{accepted} frames versus abort (abstain) rate
(\texttt{selective\_risk}). We compare it against the fixed-confidence baseline gate that accepts a frame
iff the mean detection confidence exceeds a threshold (\texttt{selective\_risk\_fixed}), at matched abort
rate, and we frame the policy as risk-controlled in the sense of distribution-free risk
control~\cite{angelopoulos2022rcps} and selective prediction with a reject option~\cite{geifman2017selective}.

\noindent\textbf{Open-loop kinematic-bicycle replay (honest, E6).} We have no field-robot odometry, so
E6 is explicitly an \emph{open-loop policy replay}, not a closed-loop field result. Over an ordered
frame sequence we feed the per-frame heading command through one kinematic-bicycle / pure-pursuit
step~\cite{rajamani2012} and accumulate a lateral-deviation surrogate (\texttt{bicycle\_replay} /
\texttt{simulate\_cross\_track}); under the abstain policy the controller holds the last trusted heading
instead of acting on an untrusted (possibly wrong-row) estimate. We report the integrated lateral
deviation and the count of large-heading-error accepted commands for the conformal policy versus a fixed
threshold; the claim is only that the abstain rule reduces large-error commands, stated as an open-loop
replay.

% =====================================================================
\subsection{Edge deployment}
\label{sec:setup-edge}
% =====================================================================

\noindent The detector is deployed on two edge platforms (Table~\ref{tab:setup-edge}): TensorRT
(FP16 and INT8) on an NVIDIA Jetson Nano~\cite{tensorrt}, and RKNN INT8 on a Rockchip
NPU~\cite{rknntoolkit}, with post-training INT8 quantization following standard integer-arithmetic
inference~\cite{jacob2018int8}. For each we report latency (mean\,$\pm$\,std and P95), throughput (FPS),
and input resolution. Because the trust layer is CPU post-NMS arithmetic over quantities the detector
already outputs, it adds no operators to the exported NPU graph, so the exported detector is a vanilla
YOLO and the on-device cost is the detector's alone; structural pruning~\cite{cth14,ct17} and efficient
backbones~\cite{ct16,ct19} are optional efficiency levers reported on the same axes if applied.
Critically, we re-fit the conformal quantile $\hat q$ \emph{after} INT8 conversion on a small
calibration set and verify that realised coverage survives quantization (E5), since INT8 perturbs the
detected centroids and hence the scores.

\begin{table}[!ht]
\centering
\caption{Edge deployment on the two named platforms. The trust layer adds no NPU operators; the
exported graph is a vanilla YOLO. \emph{Results to be completed.}}
\label{tab:setup-edge}
\resizebox{\columnwidth}{!}{\begin{tabular}{llcccc}
\hline
\textbf{Platform} & \textbf{Precision} & \textbf{Latency mean$\pm$std (ms)} & \textbf{P95 (ms)} & \textbf{FPS} & \textbf{Coverage post-INT8} \\
\hline
Jetson Nano (TensorRT) & FP16 & -- & -- & -- & -- \\
Jetson Nano (TensorRT) & INT8 & -- & -- & -- & -- \\
Rockchip (RKNN) & INT8 & -- & -- & -- & -- \\
\hline
\end{tabular}}
\end{table}

% =====================================================================
\subsection{Evaluation protocol, metrics, and statistics}
\label{sec:setup-metrics}
% =====================================================================

\noindent\textbf{Metrics.} The primary metric is heading error in degrees (controller-independent).
Line-fit RMSE is reported in pixels (and cm where the homography of Section~\ref{sec:setup-homography}
permits), evaluated \emph{including} the look-ahead row. Cross-track is a static geometric proxy: the
lateral offset between the predicted and GT lines at the look-ahead row (px, and cm where calibrated).
For the conformal layer we report realised vs.\ nominal coverage at $\alpha\in\{0.10,0.05\}$ and interval
\emph{width} (sharpness, deg) at matched coverage. For the policy we report heading error on accepted
frames and the large-heading-error rate ($>5^\circ$) versus abort rate. Detection mAP@50 and mAP@50--95
are secondary diagnostics; edge latency/throughput are reported on the named devices
(Section~\ref{sec:setup-edge}).

\noindent\textbf{Statistical protocol.} The evaluation script (\texttt{eval\_guidance.py}) writes one
CSV row per frame so that the shared-checkpoint readout arms are compared \emph{paired}, and
\texttt{eval\_conformal.py} runs the conformal experiments on those CSVs. Point-accuracy comparisons
(E1) use $\geq 5$ seeds; because $n$ is small for normality, significance uses the paired
\emph{Wilcoxon signed-rank} test with \emph{Holm--Bonferroni} correction across the confirmatory family,
and we report \emph{Cliff's $\delta$} effect sizes with 95\% bootstrap confidence intervals, not merely
$p<0.05$. Coverage proportions are reported with \emph{Clopper--Pearson} (and Wilson) confidence
intervals. The calibration/test split is by base image on CRDLD and by clip on the field video to
respect exchangeability; any claim whose $\Delta$ is smaller than its standard deviation is downgraded in
wording.

\noindent\textbf{Experiment map (E0--E6).} The protocol is organised into seven experiments, reported in
Section~\ref{sec:results}: \textbf{E0} trust-score informativeness --- Spearman $\rho(s,\text{realised
heading error})>0$, with $\sigma^2_{\mathrm{DFL}}$ ($\rho=-0.13$) as the negative control
(Section~\ref{sec:res-e0-trustinfo}); \textbf{E1} point-accuracy floor --- heading error for
equalLS/ransac/irls(deployed)/calm/qual/oracle on CRDLD test, $\geq 5$ seeds, paired Wilcoxon $+$ Holm
$+$ Cliff's $\delta$ (Section~\ref{sec:res-e1-accuracy}); \textbf{E2} coverage --- realised vs.\ nominal
at $\alpha\in\{0.10,0.05\}$, Clopper--Pearson CIs, marginal and by difficulty, split by base image
(Section~\ref{sec:res-e2-coverage}); \textbf{E3} sharpness / score-ablation --- interval width at matched
coverage for $s$ vs.\ residual vs.\ $\mathrm{conf}^2$ vs.\ $\sigma_{\mathrm{DFL}}$ vs.\
RANSAC-consensus, where the proposed score should be tightest (Section~\ref{sec:res-e3-sharpness});
\textbf{E4} selective-risk --- error on accepted frames vs.\ abort rate, conformal-width gate vs.\ the
fixed-confidence gate at matched abort (Section~\ref{sec:res-e4-selrisk}); \textbf{E5} deployment ---
Jetson/TensorRT and RKNN INT8 latency/FPS and coverage-survives-INT8 (Section~\ref{sec:res-e5-deploy});
and \textbf{E6} open-loop kinematic-bicycle policy replay (Section~\ref{sec:res-e6-replay}).
